\documentclass[border=12pt]{standalone}
\usepackage{amsmath}
\usepackage{circuitikz}
\ctikzset{bipoles/length=1.4cm}
\begin{document}
\begin{circuitikz}[american, line width=0.9pt]
  %==== CASO +M : puntos enfrentados (ambos arriba) ====
  % primario: i1 entra por el terminal con punto (arriba)
  \draw (0,3.7) node[ocirc]{} -- (0,3.0);
  \draw[-{Latex[length=2.6mm]}] (0,3.55) -- (0,3.05); \node[left=1pt] at (0,3.35){$i_1$};
  \draw (0,3.0) to[L, l_=$L_1$] (0,0.5) -- (0,-0.1) node[ocirc]{};
  \fill (0.22,2.65) circle(2.3pt);
  % secundario abierto: v2 con + arriba
  \draw (2,3.7) node[ocirc]{} -- (2,3.0); \node[right=2pt] at (2.0,3.45){$+$};
  \draw (2,3.0) to[L, l=$L_2$] (2,0.5) -- (2,-0.1) node[ocirc]{}; \node[right=2pt] at (2.0,0.05){$-$};
  \fill (1.78,2.65) circle(2.3pt);
  \draw[-{Latex[length=2.4mm]}] (2.75,3.3) -- (2.75,0.3); \node[right=1pt] at (2.75,1.8){$v_2$};
  \node[align=center] at (1.4,-1.05){\small puntos enfrentados\\[1pt]$v_2=+M\,\dfrac{di_1}{dt}$};
  %==== CASO -M : puntos opuestos ====
  \begin{scope}[xshift=7cm]
    \draw (0,3.7) node[ocirc]{} -- (0,3.0);
    \draw[-{Latex[length=2.6mm]}] (0,3.55) -- (0,3.05); \node[left=1pt] at (0,3.35){$i_1$};
    \draw (0,3.0) to[L, l_=$L_1$] (0,0.5) -- (0,-0.1) node[ocirc]{};
    \fill (0.22,2.65) circle(2.3pt);
    \draw (2,3.7) node[ocirc]{} -- (2,3.0); \node[right=2pt] at (2.0,3.45){$+$};
    \draw (2,3.0) to[L, l=$L_2$] (2,0.5) -- (2,-0.1) node[ocirc]{}; \node[right=2pt] at (2.0,0.05){$-$};
    \fill (1.78,0.85) circle(2.3pt);
    \draw[-{Latex[length=2.4mm]}] (2.75,3.3) -- (2.75,0.3); \node[right=1pt] at (2.75,1.8){$v_2$};
    \node[align=center] at (1.4,-1.05){\small puntos opuestos\\[1pt]$v_2=-M\,\dfrac{di_1}{dt}$};
  \end{scope}
\end{circuitikz}
\end{document}
